% GENERATED FILE. Edit canonical lessons, then run npm run generate:books.
% canonical-source-hash: fnv1a64:706bcd04988af00b
% canonical-lessons: TE-C166-konasagu, TE-C166-patincu, TE-C166-diddu, TE-C166-gonugu, TE-C166-vidipovu

\chapter{To continue, To observe, To correct, To mutter, To separate}
\label{ch:te-to-continue-to-observe-to-correct-to-mutter-to-separate}
\hlchaptermodality{166}

\begin{chapteropening}
\textbf{By the end of this chapter:} \emph{I can say to continue, to observe, to correct, to mutter and to separate.}

Complete the last lesson of chapter 166: I can say to continue, to observe, to correct, to mutter and to separate.
\end{chapteropening}

\section[\texorpdfstring{\te{కొనసాగు} (konasāgu)}{కొనసాగు (konasāgu)}]{\te{కొనసాగు} (konasāgu) — to continue}

\label{lesson:TE-C166-konasagu}



\begin{quote}
Before the new one: say the Telugu for to prick, then the Telugu for to knit.
\end{quote}

\subsection*{You'll want to know: \te{కొనసాగు}}
\textbf{\te{కొనసాగు}} — \emph{konasāgu} — \textquotedblleft{}to continue\textquotedblright{}.

\begin{grammarlens}[title={What you've built}]
One more word for this chapter.
\end{grammarlens}

\subsection*{Your turn}
\begin{itemize}
\raggedright
  \item \emph{Say it:} \emph{konasāgu}
  \item \emph{Say it:} \emph{konasāgu}, once more
  \item \emph{Say it:} say \emph{konasāgu}
  \item \emph{Recall:} say the Telugu for to prick, then the Telugu for to knit, then say \emph{konasāgu} again
\end{itemize}

\begin{tcolorbox}[breakable,colback=teal!4,colframe=teal!35!black,title={Before you move on}]
What does \te{కొనసాగు} mean? (\textquotedblleft{}To continue\textquotedblright{}.) Say it once more.
\end{tcolorbox}
\section[\texorpdfstring{\te{పాటించు} (pāṭiñcu)}{పాటించు (pāṭiñcu)}]{\te{పాటించు} (pāṭiñcu) — to observe, to follow (rules)}

\label{lesson:TE-C166-patincu}



\begin{quote}
Before the new one: say the Telugu for to knit, then the Telugu for to continue.
\end{quote}

\subsection*{You'll want to know: \te{పాటించు}}
\textbf{\te{పాటించు}} — \emph{pāṭiñcu} — \textquotedblleft{}to observe, to follow (rules)\textquotedblright{}.

\begin{grammarlens}[title={What you've built}]
One more word for this chapter.
\end{grammarlens}

\subsection*{Your turn}
\begin{itemize}
\raggedright
  \item \emph{Say it:} \emph{pāṭiñcu}
  \item \emph{Say it:} \emph{pāṭiñcu}, once more
  \item \emph{Say it:} say \emph{pāṭiñcu}
  \item \emph{Recall:} say the Telugu for to knit, then the Telugu for to continue, then say \emph{pāṭiñcu} again
\end{itemize}

\begin{tcolorbox}[breakable,colback=teal!4,colframe=teal!35!black,title={Before you move on}]
What does \te{పాటించు} mean? (\textquotedblleft{}To observe, to follow (rules)\textquotedblright{}.) Say it once more.
\end{tcolorbox}
\section[\texorpdfstring{\te{దిద్దు} (diddu)}{దిద్దు (diddu)}]{\te{దిద్దు} (diddu) — to correct (written work)}

\label{lesson:TE-C166-diddu}



\begin{quote}
Before the new one: say the Telugu for to continue, then the Telugu for to observe.
\end{quote}

\subsection*{You'll want to know: \te{దిద్దు}}
\textbf{\te{దిద్దు}} — \emph{diddu} — \textquotedblleft{}to correct (written work)\textquotedblright{}.

\begin{grammarlens}[title={What you've built}]
One more word for this chapter.
\end{grammarlens}

\subsection*{Your turn}
\begin{itemize}
\raggedright
  \item \emph{Say it:} \emph{diddu}
  \item \emph{Say it:} \emph{diddu}, once more
  \item \emph{Say it:} say \emph{diddu}
  \item \emph{Recall:} say the Telugu for to continue, then the Telugu for to observe, then say \emph{diddu} again
\end{itemize}

\begin{tcolorbox}[breakable,colback=teal!4,colframe=teal!35!black,title={Before you move on}]
What does \te{దిద్దు} mean? (\textquotedblleft{}To correct (written work)\textquotedblright{}.) Say it once more.
\end{tcolorbox}
\section[\texorpdfstring{\te{గొణుగు} (goṇugu)}{గొణుగు (goṇugu)}]{\te{గొణుగు} (goṇugu) — to mutter, to grumble}

\label{lesson:TE-C166-gonugu}



\begin{quote}
Before the new one: say the Telugu for to observe, then the Telugu for to correct.
\end{quote}

\subsection*{You'll want to know: \te{గొణుగు}}
\textbf{\te{గొణుగు}} — \emph{goṇugu} — \textquotedblleft{}to mutter, to grumble\textquotedblright{}.

\begin{grammarlens}[title={What you've built}]
One more word for this chapter.
\end{grammarlens}

\subsection*{Your turn}
\begin{itemize}
\raggedright
  \item \emph{Say it:} \emph{goṇugu}
  \item \emph{Say it:} \emph{goṇugu}, once more
  \item \emph{Say it:} say \emph{goṇugu}
  \item \emph{Recall:} say the Telugu for to observe, then the Telugu for to correct, then say \emph{goṇugu} again
\end{itemize}

\begin{tcolorbox}[breakable,colback=teal!4,colframe=teal!35!black,title={Before you move on}]
What does \te{గొణుగు} mean? (\textquotedblleft{}To mutter, to grumble\textquotedblright{}.) Say it once more.
\end{tcolorbox}
\section[\texorpdfstring{\te{విడిపోవు} (viḍipōvu)}{విడిపోవు (viḍipōvu)}]{\te{విడిపోవు} (viḍipōvu) — to separate, to part}

\label{lesson:TE-C166-vidipovu}



\begin{quote}
Before the new one: say the Telugu for to correct, then the Telugu for to mutter.
\end{quote}

\subsection*{You'll want to know: \te{విడిపోవు}}
\textbf{\te{విడిపోవు}} — \emph{viḍipōvu} — \textquotedblleft{}to separate, to part\textquotedblright{}.

\begin{grammarlens}[title={What you've built}]
One more word for this chapter.
\end{grammarlens}

\subsection*{Your turn}
\begin{itemize}
\raggedright
  \item \emph{Say it:} \emph{viḍipōvu}
  \item \emph{Say it:} \emph{viḍipōvu}, once more
  \item \emph{Say it:} say \emph{viḍipōvu}
  \item \emph{Recall:} say the Telugu for to correct, then the Telugu for to mutter, then say \emph{viḍipōvu} again
\end{itemize}

\begin{tcolorbox}[breakable,colback=teal!4,colframe=teal!35!black,title={Before you move on}]
What does \te{విడిపోవు} mean? (\textquotedblleft{}To separate, to part\textquotedblright{}.) Say it once more.
\end{tcolorbox}
